\usepackage[dvipsnames]{xcolor}
\usepackage{amsmath}   % 数学公式
\usepackage{amsthm}    % 定理环境
\usepackage{amssymb}   % 更多公式符号
\usepackage{graphicx}  
\usepackage{wrapfig}   % 插图
\usepackage{mathrsfs}  % 数学字体
\usepackage{physics}   % 物理符号
\usepackage{booktabs}
\usepackage{tabularx}
\usepackage{longtable} % 支持跨页表格
\usepackage{array}     % 增强表格功能
\usepackage{multirow}
\usepackage{makecell}
\usepackage{tikz}
\usepackage{tikz-3dplot} % 引入 tikz-3dplot 宏包
\usepackage{subcaption} % 用于子图排版
\usetikzlibrary{arrows.meta, positioning, shapes.geometric}
\usepackage{enumitem}  % 列表
\usepackage{geometry}  % 页面调整
\usepackage{ragged2e}
\usepackage{unicode-math}
\usepackage[hidelinks,
colorlinks=true,
allcolors=blue,
pdfstartview=Fit,
breaklinks=true]{hyperref}

\allowdisplaybreaks % 允许公式跨页换行

\captionsetup{hypcap=false} % 禁用 hypcap 功能，消除\captionof 的警告

\graphicspath{ {flg/},{../flg/}, {config/}, {../config/} }  % 配置图形文件检索目录
\linespread{1.2} % 行高

% 页面设置
\geometry{
    top=15mm,
    bottom=15mm,
    left=15mm,
    right=15mm,
    headheight=4mm,
    headsep=4mm,
    footskip=8mm
}

% 调整章节标题格式
\ctexset{
    chapter = {
        format      = \huge\bfseries\centering,  % 字体大小 + 加粗 + 居中
        name={,、},
		break = {}, 
		beforeskip = 20pt,                       % 标题前的垂直间距
        afterskip  = 20pt                        % 标题后的垂直间距
    },
    section = {
        format     = \Large\bfseries\centering 
    },
    subsubsection= {
        format     = \large\bfseries\raggedright, % 字体大小 + 加粗 + 左对齐
    },
    paragraph = {
        format     = \normalsize\bfseries\raggedright,
        runin      = false
    }
}

% 设置列表环境的上下间距
\setenumerate[1]{itemsep=5pt,partopsep=0pt,parsep=\parskip,topsep=5pt}
\setitemize[1]{itemsep=5pt,partopsep=0pt,parsep=\parskip,topsep=5pt}
\setdescription{itemsep=5pt,partopsep=0pt,parsep=\parskip,topsep=5pt}


% ########## 定理环境 start ####################################
\theoremstyle{definition}
\newtheorem{definition}{\indent 定义}[section]

\newtheorem{add}{\indent 补充}[section]

\newtheorem{lemma}{\indent 引理}[section]    % 引理 定理 推论 准则 共用一个编号计数
\newtheorem{theorem}[lemma]{\indent 定理}
\newtheorem{corollary}[lemma]{\indent 推论}
\newtheorem{criterion}[lemma]{\indent 准则}
\newtheorem{law}[lemma]{\indent 定律}
\newtheorem{principle}[lemma]{\indent 原理}
\newtheorem{proposition}[lemma]{\indent 命题}

\newtheorem{example}{\indent \textcolor{SeaGreen}{例}}[section]

% 无编号定理环境
\newenvironment{note}{\par\textbf{注：}}{}
\newenvironment{explain}{\par\textbf{解释：}}{}
\renewenvironment{proof}{\par\textbf{证明：}}{\qed\par}
\newenvironment{solution}{\par\textbf{解：}}{\qed\par}
\newenvironment{known}{\par\textbf{已知：}}{}
\renewenvironment{derivative}{\par\textbf{推导：}}{\qed\par}

% ######### 定理环境 end  #####################################

% ↓↓↓↓↓↓↓↓↓↓↓↓↓↓↓↓↓ 以下是自定义的命令  ↓↓↓↓↓↓↓↓↓↓↓↓↓↓↓↓

% 用于调整表格的高度  使用 \hline\xrowht{25pt}
\newcommand{\xrowht}[2][0]{\addstackgap[.5\dimexpr#2\relax]{\vphantom{#1}}}

% 表格环境内长内容换行
\newcommand{\tabincell}[2]{\begin{tabular}{@{}#1@{}}#2\end{tabular}}

% 使用\linespread{1.5} 之后 cases 环境的行高也会改变，重新定义一个 ca 环境可以自动控制 cases 环境行高
% \newenvironment{ca}[1][1]{\linespread{#1} \selectfont \begin{cases}}{\end{cases}}
% 和上面一样
% \newenvironment{vx}[1][1]{\linespread{#1} \selectfont \begin{vmatrix}}{\end{vmatrix}}

\def\d{\textup{d}} % 直立体 d 用于微分符号 dx
\def\R{\mathbb{R}} % 实数域

% 数学 平行 符号
\newcommand{\pll}{\kern 0.56em/\kern -0.8em /\kern 0.56em}
\newcommand{\pa}{\partial}
\newcommand{\vvec}[1]{\symbf{#1}}
\newcommand{\ten}[1]{\mathbb{#1}}
% \newcommand{\ten}{\overset{\rightharpoonup\!\!\!\! \rightharpoonup}}
% 中文平行 parallel
\newcommand*\zhparallel{%
  \mathrel{\text{\tikz[baseline,line cap=round] \draw (0em,-0.3ex) -- (.4em,1.7ex) (.2em,-0.3ex) -- (.6em,1.7ex);}}%
}
\newcommand{\X}{\mathtt{X}}
\newcommand{\mT}{\raisebox{0.1ex}{$\scriptstyle -T$}} % 调整高度为 0.1ex
\newcommand{\lmT}{\raisebox{-0.85ex}{$\scriptstyle -T$}} % 调整高度为 -0.85ex
\newcommand{\mone}{\raisebox{0.1ex}{$\scriptstyle -1$}} % 调整高度为 0.1ex
\newcommand{\lmone}{\raisebox{-0.85ex}{$\scriptstyle -1$}} % 调整高度为 -0.85ex
\newcommand{\lsup}[1]{\raisebox{-0.1ex}{$\scriptstyle #1$}}
\newcommand{\llsup}[1]{\raisebox{-0.85ex}{$\scriptstyle #1$}}
\newcommand{\mathminus}{\!\!-\!\!} % 数学环境连字符

%颜色
\definecolor{b1}{RGB}{0,191,255}    %deep sky blue
\definecolor{g1}{RGB}{120,163,85}    %柳色



% 定义双曲函数及其反函数
\DeclareMathOperator{\arsh}{arsinh} % 反双曲正弦
\DeclareMathOperator{\arch}{arcosh} % 反双曲余弦
\DeclareMathOperator{\arth}{artanh} % 反双曲正切
\DeclareMathOperator{\arcoth}{arcoth} % 反双曲余切
